
\documentclass{article}
\usepackage{colortbl}
\usepackage{makecell}
\usepackage{multirow}
\usepackage{supertabular}

\begin{document}

\newcounter{utterance}

\centering \large Interaction Transcript for game `bbh', experiment `navigate', episode 42 with qwen3.5{-}2b.
\vspace{24pt}

{ \footnotesize  \setcounter{utterance}{1}
\setlength{\tabcolsep}{0pt}
\begin{supertabular}{c@{$\;$}|p{.15\linewidth}@{}p{.15\linewidth}p{.15\linewidth}p{.15\linewidth}p{.15\linewidth}p{.15\linewidth}}
    \# & \multicolumn{2}{c}{Player} && \multicolumn{2}{c}{Game Master} \\
    \hline

    \theutterance \stepcounter{utterance}  
    & & & \multicolumn{4}{p{0.6\linewidth}}{
        \cellcolor[rgb]{0.9,0.9,0.9}{
            \makecell[{{p{\linewidth}}}]{
                \texttt{\tiny{[P1$\langle$GM]}}
                \texttt{Given a series of navigation instructions, determine whether one would end up back at the starting point.} \\
\\ 
\texttt{Q: If you follow these instructions, do you return to the starting point? Turn left. Turn around. Turn left. Take 7 steps. Take 2 steps. Take 4 steps. Take 8 steps.} \\
\texttt{Options:} \\
\texttt{{-} Yes} \\
\texttt{{-} No} \\
\texttt{A: No} \\
\\ 
\texttt{Q: If you follow these instructions, do you return to the starting point? Turn around. Take 1 step. Take 6 steps. Turn around. Take 6 steps. Take 9 steps. Take 1 step.} \\
\texttt{Options:} \\
\texttt{{-} Yes} \\
\texttt{{-} No} \\
\texttt{A: No} \\
\\ 
\texttt{Q: If you follow these instructions, do you return to the starting point? Always face forward. Take 2 steps right. Take 9 steps left. Take 7 steps right.} \\
\texttt{Options:} \\
\texttt{{-} Yes} \\
\texttt{{-} No} \\
\texttt{A: Yes} \\
\\ 
\texttt{Q: If you follow these instructions, do you return to the starting point? Always face forward. Take 7 steps left. Take 2 steps backward. Take 7 steps backward. Take 7 steps backward. Take 3 steps forward.} \\
\texttt{Options:} \\
\texttt{{-} Yes} \\
\texttt{{-} No} \\
\texttt{A:} \\
            }
        }
    }
     \\ \\

    \theutterance \stepcounter{utterance}  
    & \multicolumn{4}{p{0.6\linewidth}}{
        \cellcolor[rgb]{0.561,0.737,0.561}{
            \makecell[{{p{\linewidth}}}]{
                \texttt{\tiny{[P1$\rangle$GM]}}
                \texttt{To determine if you return to the starting point, let's analyze the movements relative to the starting position (0):} \\
\\ 
\texttt{1.  **Always face forward**: This sets the direction for the subsequent steps.} \\
\texttt{2.  **Take 7 steps left**: You move 7 units in the negative direction.} \\
\texttt{    *   Current position: \${-}7\$} \\
\texttt{3.  **Take 2 steps backward**: "Backward" is the opposite of "forward" (or "left" if facing forward). So, you move 2 units in the negative direction.} \\
\texttt{    *   Current position: \${-}7 {-} 2 = {-}9\$} \\
\texttt{4.  **Take 7 steps backward**: You move 7 units in the negative direction.} \\
\texttt{    *   Current position: \${-}9 {-} 7 = {-}16\$} \\
\texttt{5.  **Take 7 steps backward**: You move 7 units in the negative direction.} \\
\texttt{    *   Current position: \${-}16 {-} 7 = {-}23\$} \\
\texttt{6.  **Take 3 steps forward**: "Forward" is the opposite of "backward". So, you move 3 units in the positive direction.} \\
\texttt{    *   Current position: \${-}23 + 3 = {-}20\$} \\
\\ 
\texttt{The final position (\${-}20\$) is not the same as the starting position (\$0\$).} \\
\\ 
\texttt{A: No} \\
            }
        }
    }
    & & \\ \\

    \theutterance \stepcounter{utterance}  
    & & & \multicolumn{2}{p{0.3\linewidth}}{
        \cellcolor[rgb]{0.9,0.9,0.9}{
            \makecell[{{p{\linewidth}}}]{
                \texttt{\tiny{[GM$|$GM]}}
                \texttt{Target: No} \\
            }
        }
    }
    & & \\ \\

    \theutterance \stepcounter{utterance}  
    & & & \multicolumn{2}{p{0.3\linewidth}}{
        \cellcolor[rgb]{0.9,0.9,0.9}{
            \makecell[{{p{\linewidth}}}]{
                \texttt{\tiny{[GM$|$GM]}}
                \texttt{game\_result = LOSE} \\
            }
        }
    }
    & & \\ \\

\end{supertabular}
}

\end{document}
